% Diapositiva 16

% Luchando los dias
%============================================
% 15. Explicar procedimiento de las amplitudes ideales
%=============================================

\begin{frame}{Cálculo de Amplitudes Ideales ($A_{ideal}$)}

    \begin{columns}[T]
        %========================================
        % COLUMNA IZQUIERDA: Tabla de Datos Ri
        %========================================
        \column{0.48\textwidth}
        \begin{block}{Entradas: Coordenadas y Distancia}
            \centering
            \scriptsize 
            \begin{tabular}{@{}lcccc@{}}
                \toprule
                \textbf{Sens.} & \textbf{$x_i$} & \textbf{$y_i$} & \textbf{$z_i$} & \textbf{$R_i$} \\ \midrule
                S1  & -4.0 & -3.0 &  0.2 & 5.7940 \\
                S2  & -2.5 &  2.8 & -0.1 & 5.2583 \\
                S3  &  0.0 &  4.0 &  0.1 & 5.0912 \\
                S4  &  3.5 &  3.0 & -0.2 & 4.5321 \\
                S5  &  4.5 & -1.5 &  0.0 & 3.5482 \\
                S6  &  2.0 & -4.0 &  0.3 & 3.5833 \\
                S7  & -1.0 & -4.5 & -0.1 & 4.4193 \\
                S8  & -4.5 &  1.0 &  0.2 & 6.1172 \\
                S9  &  0.5 &  0.5 &  0.0 & 1.8412 \\
                S10 &  2.8 &  0.8 & -0.3 & 2.4000 \\
                S11 & -1.8 & -1.2 &  0.1 & 3.2558 \\
                S12 &  1.0 &  3.2 &  0.2 & 4.2107 \\ \bottomrule
            \end{tabular}
        \end{block}

        %========================================
        % COLUMNA DERECHA: Procedimiento Paso a Paso
        %========================================
        \column{0.48\textwidth}
        \begin{exampleblock}{Procedimiento ($S_1$)}
            \centering
            \footnotesize
            % Se redujo el salto de línea de 0.3cm a 0.1cm para ahorrar espacio vertical
            \begin{align*}
                A_{ideal,i} &= A_0 \frac{e^{-R_i}}{R_i} \\[0.1cm]
                A_{ideal,1} &= 10 \frac{e^{-5.7940}}{5.7940} \\[0.1cm]
                A_{ideal,1} &\approx 10 \frac{0.003045}{5.7940} \\[0.1cm]
                \mathbf{A_{ideal,1}} &\approx \mathbf{0.005257}
            \end{align*}
        \end{exampleblock}

        \vspace{0.1cm} % Reducido para acercar la nota al bloque superior

        \begin{block}{Nota del Modelo}
            \scriptsize
            Cálculo determinista basado en el decaimiento exponencial y dispersión geométrica.
        \end{block}
    \end{columns}

\end{frame}